\begin{tabular}{llrrrrr}
\toprule
\textbf{Det.} & \textbf{Fuente} & \textbf{$n_V$} & \textbf{$b$} & \textbf{$\chi^2$ potencia} & \textbf{$\chi^2$ recta} & \textbf{Peor residuo [\%]} \\
\midrule
A & Co-60 & 4 & $9.37 \pm 0.03$ & 53 & 1707 & 1.1 \\
A & Cs-137 & 3 & $10.42 \pm 0.04$ & 0 & 358 & 0.1 \\
B & Co-60 & 3 & $9.44 \pm 0.04$ & 5 & 244 & 0.3 \\
B & Cs-137 & 3 & $9.56 \pm 0.05$ & 0 & 192 & 0.1 \\
C & Co-60 & 3 & $8.70 \pm 0.05$ & 168 & 638 & 2.1 \\
C & Cs-137 & 3 & $9.34 \pm 0.05$ & 4 & 245 & 0.3 \\
D & Co-60 & 3 & $10.10 \pm 0.04$ & 493 & 1749 & 2.9 \\
D & Cs-137 & 3 & $7.73 \pm 0.05$ & 22 & 31 & 0.6 \\
\bottomrule
\end{tabular}
\par\smallskip
\begin{minipage}{0.95\linewidth}\footnotesize Ajuste de $D(V)$ con las dos hipotesis. El $\chi^2$ de la recta es una a dos ordenes de magnitud mayor en todos los casos, o sea que la ley de potencia gana sin ambiguedad. Un PMT tipico tiene $b$ entre 6 y 10.\end{minipage}
